% !Mode:: "TeX:UTF-8"

\chapter*{\centering \sanhao\song\bfseries 摘\quad 要}

\xiaosi\song
\setstretch{1.5}

本文围绕“电动轮椅能否在狭窄楼梯转角处完成 90 度转弯”这一空间推理任务，构建多层级 Prompt 数据集，将大语言模型的回答视为可挖掘的行为数据，分析模型在自然语言描述、参数化描述和数学建模描述之间的判断稳定性与认知边界。实验设计包含 20 个空间场景、3 个描述层级和 5 类无关干扰条件，形成 300 条基础 Prompt；随后采集模型回复，并抽取最终判断、推理步骤、公式使用、坐标系使用和错误类型等行为特征。

在数据挖掘阶段，本文设计了准确率统计、层级一致性、噪声翻转率、决策树、K-Means 聚类和 FP-Growth 关联规则等分析方法，用于识别模型在空间建模任务中的典型失效模式。该框架能够回答三个核心问题：第一，形式化几何描述是否一定提升模型判断准确率；第二，材质颜色、环境噪声、光照和情绪压力等无关信息是否会诱发模型判断翻转；第三，不同模型在空间推理任务中呈现出怎样的行为边界和错误模式。

当前报告基于项目已有代码和数据产物搭建完整分析框架，后续可在真实模型多轮采集完成后补充具体统计数值、典型错误样例和最终结论。该项目展示了如何将大语言模型回复转化为结构化行为数据，并使用数据挖掘方法分析模型能力边界。

\vspace{0.8cm}

\noindent {\hei\xiaosi 关键词：} 大语言模型；空间推理；数据挖掘；Prompt 工程；关联规则
